\chapter{Complexity — Circuit Sat}
%%% generated by the blueprint pipeline from TCSlib.Complexity.CircuitComplexity.CircuitSat
\section{Overview}

Declarations of the Lean module \texttt{TCSlib.Complexity.CircuitComplexity.CircuitSat}.

\section{Declarations}

\begin{definition}[Tree-circuit satisfiability]
\lean{BoolCircuit.TreeCircuit.Satisfiable}
\label{BoolCircuit.TreeCircuit.Satisfiable}
\leanok
\uses{BoolCircuit.TreeCircuit, BoolCircuit.TreeCircuit.eval}
A Boolean circuit tree on $n$ variables is satisfiable if there exists an assignment of
Boolean values to the $n$ variables for which the circuit evaluates to true.
\end{definition}

\begin{definition}[Satisfiable circuits]
\lean{BoolCircuit.CktSat}
\statementsource{Vad12}{
  pdf-bf2ce5422088-p029-b001,
  pdf-bf2ce5422088-p031-b004
}
\label{BoolCircuit.CktSat}
\leanok
\uses{BoolCircuit.TreeCircuit, BoolCircuit.TreeCircuit.Satisfiable}
For each natural number $n$, let the set of satisfiable circuits consist of the pairs
$(n,C)$ where $C$ is a Boolean circuit on $n$ variables for which some assignment of
truth values to the $n$ inputs makes $C$ evaluate to true.
\end{definition}

\begin{theorem}[Circuit Satisfiability Membership]
\lean{BoolCircuit.mem_cktSat_iff}
\label{BoolCircuit.mem_cktSat_iff}
\leanok
\uses{BoolCircuit.CktSat, BoolCircuit.TreeCircuit, BoolCircuit.TreeCircuit.Satisfiable}
\difficulty{1}
For a Boolean circuit tree on $n$ variables, membership of the corresponding
variable-count and circuit pair in the set of satisfiable circuit trees is equivalent to
the circuit tree being satisfiable.
\end{theorem}

\begin{definition}[Circuit variables]
\lean{BoolCircuit.CktVar}
\label{BoolCircuit.CktVar}
\leanok
\uses{BoolCircuit.TreeCircuit}
For each natural number $n$, a circuit variable is either an input indexed by one of the
$n$ input positions or a gate labelled by a Boolean circuit on those inputs.
\end{definition}

\begin{definition}[Positive circuit literal]
\lean{BoolCircuit.Lit.toCktLiteral}
\label{BoolCircuit.Lit.toCktLiteral}
\leanok
\uses{BoolCircuit.CktVar, BoolCircuit.Lit, SATTo3SAT.Literal}
For a circuit literal $l$ on $n$ variables, its associated literal over circuit
variables is the positive occurrence of input $l_{\mathrm{idx}}$ if $l$ is positive, and
the negative occurrence of that input if $l$ is negated.
\end{definition}

\begin{definition}[Negated circuit literal]
\lean{BoolCircuit.Lit.toCktLiteralNeg}
\label{BoolCircuit.Lit.toCktLiteralNeg}
\leanok
\uses{BoolCircuit.CktVar, BoolCircuit.Lit, SATTo3SAT.Literal}
For a circuit literal $l$ on $n$ variables, its associated literal over circuit
variables is the negative occurrence of input $l_{\mathrm{idx}}$ if $l$ is positive, and
the positive occurrence of that input if $l$ is negated.
\end{definition}

\begin{definition}[Clauses for a gate]
\lean{BoolCircuit.gateClauses}
\label{BoolCircuit.gateClauses}
\leanok
\uses{BoolCircuit.CktVar, BoolCircuit.TreeCircuit, SATTo3SAT.CNFFormula}
For a gate type and a list of child circuits on $n$ variables, the gate clauses are
defined as follows. For an AND gate with output variable $g$ and child variables
$c_1,\ldots,c_k$, they comprise $(g\lor\neg c_1\lor\cdots\lor\neg c_k)$ and, for each
child $c_i$, the clause $(\neg g\lor c_i)$. For an OR gate, they comprise $(\neg g\lor
c_1\lor\cdots\lor c_k)$ and, for each child $c_i$, the clause $(g\lor\neg c_i)$.
\end{definition}

\begin{definition}[Tseitin clauses for a circuit]
\lean{BoolCircuit.TreeCircuit.tseitin}
\label{BoolCircuit.TreeCircuit.tseitin}
\leanok
\uses{BoolCircuit.CktVar, BoolCircuit.Lit.toCktLiteral, BoolCircuit.Lit.toCktLiteralNeg, BoolCircuit.TreeCircuit, BoolCircuit.gateClauses, SATTo3SAT.CNFFormula}
For a circuit on $n$ variables, its Tseitin clauses are defined recursively. For a leaf
carrying literal $l$, they are the two clauses expressing that the gate variable for the
leaf has the same truth value as $l$. For an internal gate, they consist of the clauses
for that gate together with the Tseitin clauses of each child circuit.
\end{definition}

\begin{definition}[CNF encoding of a circuit]
\lean{BoolCircuit.TreeCircuit.toCNF}
\label{BoolCircuit.TreeCircuit.toCNF}
\leanok
\uses{BoolCircuit.CktVar, BoolCircuit.TreeCircuit, BoolCircuit.TreeCircuit.tseitin, SATTo3SAT.CNFFormula}
For a Boolean circuit $C$ on $n$ variables, its CNF encoding consists of the unit clause
asserting that the gate variable for $C$ is true, together with the Tseitin clauses of
$C$.
\end{definition}

\begin{definition}[Circuit evaluation assignment]
\lean{BoolCircuit.tseitinAssignment}
\label{BoolCircuit.tseitinAssignment}
\leanok
\uses{BoolCircuit.CktVar, BoolCircuit.TreeCircuit.eval}
Given an assignment $u$ of truth values to the $n$ inputs, define a predicate on circuit
variables by declaring an input variable true exactly when $u$ assigns it true, and a
gate variable for a circuit $D$ true exactly when $D$ evaluates to true under $u$.
\end{definition}

\begin{theorem}[Circuit literal evaluation agrees with its encoding]
\lean{BoolCircuit.evalLiteral_toCktLiteral_iff}
\label{BoolCircuit.evalLiteral_toCktLiteral_iff}
\leanok
\uses{BoolCircuit.CktVar, BoolCircuit.Lit, BoolCircuit.Lit.eval, BoolCircuit.Lit.toCktLiteral, SATTo3SAT.Assignment, SATTo3SAT.evalLiteral}
\difficulty{4}
Let $n$ be a natural number, let $\alpha$ assign a truth proposition to every input or
gate variable of a tree circuit on $n$ inputs, and let $u$ be a Boolean assignment to
its inputs. Suppose that $u(i)$ is true exactly when $\alpha$ makes input $i$ true, for
every input $i$. Then for every literal $l$, the circuit literal associated with $l$ is
true under $\alpha$ exactly when $l$ evaluates to true under $u$.
\end{theorem}

\begin{theorem}[Negated circuit literal evaluation]
\lean{BoolCircuit.evalLiteral_toCktLiteralNeg_iff}
\label{BoolCircuit.evalLiteral_toCktLiteralNeg_iff}
\leanok
\uses{BoolCircuit.CktVar, BoolCircuit.Lit, BoolCircuit.Lit.eval, BoolCircuit.Lit.toCktLiteralNeg, SATTo3SAT.Assignment, SATTo3SAT.evalLiteral}
\difficulty{4}
Let $n$ be a natural number, let $\alpha$ assign a truth proposition to every input or
gate variable of a tree circuit on $n$ inputs, and let $u$ be a Boolean assignment to
its inputs. Suppose that $u(i)$ is true exactly when $\alpha$ makes input $i$ true, for
every input $i$. Then for every literal $l$, the negated circuit literal associated with
$l$ is true under $\alpha$ exactly when $l$ evaluates to false under $u$.
\end{theorem}

\begin{theorem}[Satisfaction of gate clauses]
\lean{BoolCircuit.gateClauses_satisfied}
\label{BoolCircuit.gateClauses_satisfied}
\leanok
\uses{BoolCircuit.TreeCircuit, BoolCircuit.TreeCircuit.eval, BoolCircuit.TreeCircuit.eval_node_false_iff, BoolCircuit.TreeCircuit.eval_node_true_iff, BoolCircuit.gateClauses, BoolCircuit.tseitinAssignment, SATTo3SAT.formulaSatisfied}
\difficulty{5}
Let $n$ be a natural number, let $u$ be an assignment of Boolean values to $n$
variables, let $b$ be a Boolean value, and let $cs$ be a finite list of Boolean circuit
trees on those variables. The clauses associated with $b$ and $cs$ are all satisfied by
the assignment that gives each input variable its value under $u$ and gives each circuit
gate the value of that circuit under $u$.
\end{theorem}

\begin{theorem}[Satisfaction of the Tseitin Encoding]
\lean{BoolCircuit.tseitin_satisfied}
\label{BoolCircuit.tseitin_satisfied}
\leanok
\uses{BoolCircuit.Lit.eval, BoolCircuit.Lit.toCktLiteral, BoolCircuit.Lit.toCktLiteralNeg, BoolCircuit.TreeCircuit, BoolCircuit.TreeCircuit.eval_lit, BoolCircuit.TreeCircuit.ind, BoolCircuit.TreeCircuit.tseitin, BoolCircuit.evalLiteral_toCktLiteralNeg_iff, BoolCircuit.evalLiteral_toCktLiteral_iff, BoolCircuit.gateClauses_satisfied, BoolCircuit.tseitinAssignment, SATTo3SAT.formulaSatisfied}
\difficulty{5}
Let $u$ be a Boolean assignment to the inputs of a Boolean circuit tree $C$. Assign each
input variable its value under $u$ and each gate variable the value obtained by
evaluating the circuit at that gate. Then every clause in the Tseitin encoding of $C$ is
satisfied by this assignment.
\end{theorem}

\begin{theorem}[Gate evaluation from satisfied gate clauses]
\lean{BoolCircuit.eval_node_iff_of_gateClauses}
\label{BoolCircuit.eval_node_iff_of_gateClauses}
\leanok
\uses{BoolCircuit.CktVar, BoolCircuit.TreeCircuit, BoolCircuit.TreeCircuit.eval, BoolCircuit.TreeCircuit.eval_node_false_iff, BoolCircuit.TreeCircuit.eval_node_true_iff, BoolCircuit.gateClauses, SATTo3SAT.Assignment, SATTo3SAT.clauseSatisfied, SATTo3SAT.evalLiteral, SATTo3SAT.formulaSatisfied}
\difficulty{6}
Let $n$ be a natural number, let $\alpha$ assign a truth proposition to every input or
gate variable of a tree circuit on $n$ inputs, and let $u$ be a Boolean assignment to
its inputs. For a gate label $b$ and a list $cs$ of child circuits, suppose all clauses
associated with that gate and child list are satisfied by $\alpha$. Suppose also that
each child circuit $c$ evaluates to true under $u$ exactly when $\alpha$ makes its gate
variable true. Then the circuit with gate label $b$ and children $cs$ evaluates to true
under $u$ exactly when $\alpha$ makes its gate variable true.
\end{theorem}

\begin{theorem}[Circuit evaluation from satisfied Tseitin clauses]
\lean{BoolCircuit.eval_iff_of_tseitin_satisfied}
\label{BoolCircuit.eval_iff_of_tseitin_satisfied}
\leanok
\uses{BoolCircuit.CktVar, BoolCircuit.Lit.eval, BoolCircuit.Lit.toCktLiteral, BoolCircuit.Lit.toCktLiteralNeg, BoolCircuit.TreeCircuit, BoolCircuit.TreeCircuit.eval, BoolCircuit.TreeCircuit.eval_lit, BoolCircuit.TreeCircuit.ind, BoolCircuit.TreeCircuit.tseitin, BoolCircuit.evalLiteral_toCktLiteralNeg_iff, BoolCircuit.evalLiteral_toCktLiteral_iff, BoolCircuit.eval_node_iff_of_gateClauses, SATTo3SAT.Assignment, SATTo3SAT.formulaSatisfied}
\difficulty{5}
Let $n$ be a natural number, let $\alpha$ assign a truth proposition to every input or
gate variable of a tree circuit on $n$ inputs, and let $u$ be a Boolean assignment to
its inputs. Suppose that $u(i)$ is true exactly when $\alpha$ makes input $i$ true, for
every input $i$. If all clauses in the Tseitin formula of a tree circuit $C$ are
satisfied by $\alpha$, then $C$ evaluates to true under $u$ exactly when $\alpha$ makes
the gate variable for $C$ true.
\end{theorem}

\begin{theorem}[Circuit satisfiability and CNF encoding]
\lean{BoolCircuit.TreeCircuit.satisfiable_iff_isSatisfiable}
\label{BoolCircuit.TreeCircuit.satisfiable_iff_isSatisfiable}
\leanok
\uses{BoolCircuit.TreeCircuit, BoolCircuit.TreeCircuit.Satisfiable, BoolCircuit.TreeCircuit.toCNF, BoolCircuit.eval_iff_of_tseitin_satisfied, BoolCircuit.tseitinAssignment, BoolCircuit.tseitin_satisfied, SATTo3SAT.clauseSatisfied, SATTo3SAT.evalLiteral, SATTo3SAT.isSatisfiable}
\difficulty{4}
Let $n$ be a natural number and let $C$ be a Boolean circuit on $n$ variables. There
exists an assignment to its inputs making $C$ evaluate to true if and only if the CNF
formula encoding $C$, with its output required to be true, has a satisfying assignment.
\end{theorem}

\begin{theorem}[Circuit Satisfiability and 3-SAT]
\lean{BoolCircuit.mem_cktSat_iff_is3Satisfiable}
\label{BoolCircuit.mem_cktSat_iff_is3Satisfiable}
\leanok
\uses{BoolCircuit.CktSat, BoolCircuit.TreeCircuit, BoolCircuit.TreeCircuit.satisfiable_iff_isSatisfiable, BoolCircuit.TreeCircuit.toCNF, SATTo3SAT.SAT_to_3SAT_equivalence, SATTo3SAT.is3Satisfiable, SATTo3SAT.to3SAT}
\difficulty{1}
For a Boolean circuit tree on $n$ variables, the circuit is satisfiable if and only if
the 3-CNF formula obtained by converting its Tseitin CNF encoding to 3-CNF is
satisfiable.
\end{theorem}

\begin{theorem}[Length bound for the Tseitin formula]
\lean{BoolCircuit.TreeCircuit.tseitin_length_lt}
\label{BoolCircuit.TreeCircuit.tseitin_length_lt}
\leanok
\uses{BoolCircuit.TreeCircuit, BoolCircuit.TreeCircuit.ind, BoolCircuit.TreeCircuit.size, BoolCircuit.TreeCircuit.tseitin, BoolCircuit.gateClauses}
\difficulty{5}
For every Boolean circuit $C$ on $n$ variables, the length of its Tseitin formula is
strictly less than $3\,\mathrm{size}(C)$.
\end{theorem}

\begin{theorem}[Clause count in the circuit CNF translation]
\lean{BoolCircuit.TreeCircuit.toCNF_length_le}
\label{BoolCircuit.TreeCircuit.toCNF_length_le}
\leanok
\uses{BoolCircuit.TreeCircuit, BoolCircuit.TreeCircuit.size, BoolCircuit.TreeCircuit.toCNF, BoolCircuit.TreeCircuit.tseitin_length_lt}
\difficulty{2}
For every Boolean circuit $C$ on $n$ variables, the length of its conjunctive normal
form translation is at most three times the size of $C$.
\end{theorem}
